\documentclass[12pt, a4paper]{article}
\usepackage[utf8]{inputenc}
\usepackage[T1]{fontenc}
\usepackage{amsmath, amssymb, amsthm}
\usepackage{geometry}
\geometry{margin=2.5cm}
\usepackage{hyperref}
\usepackage{listings}
\usepackage{booktabs}
\usepackage{xeCJK}
\setCJKmainfont{cwTeX Q Ming}

\lstset{
  basicstyle=\ttfamily\small,
  breaklines=true,
  frame=single,
  language=R
}

\title{如何使用 Claude Code 進行經濟學研究}
\author{Sheng-Yuan Chen}
\date{2026-04-07}

\begin{document}
\maketitle

\section*{前言}

經濟學研究的日常充滿繁瑣的資料處理工作：從各種 API 抓資料、清理格式、合併 panel、寫回歸、除錯——這些步驟耗時，卻不一定是真正產生學術價值的地方。

\textbf{Claude Code} 是 Anthropic 推出的 AI 程式助理，可以直接在終端機或 VS Code 裡運作，讀取、撰寫和執行程式碼。本文以一個真實的研究專案（台灣同婚合法化 $\times$ 創新產出：合成控制法分析）為例，示範如何用 Claude Code 加速資料整理流程。

\section{安裝與基本設定}

\subsection{安裝 Claude Code}

Claude Code 透過 npm 安裝，需要 Node.js 18+：

\begin{lstlisting}[language=bash]
# 確認 Node.js 版本
node --version

# 安裝 Claude Code CLI
npm install -g @anthropic-ai/claude-code

# 啟動（在專案資料夾內執行）
claude
\end{lstlisting}

\subsection{設定專案記憶（CLAUDE.md）}

在專案根目錄建立 \texttt{CLAUDE.md}，告訴 Claude Code 這個專案的基本規則：

\begin{lstlisting}[language=bash]
# 專案規則
- 語言：R，使用 tidyverse 風格
- 路徑：用 here::here()，不用 setwd()
- 種子：每個 script 開頭 set.seed(42)
- 輸出：圖片存 output/figures/，處理後資料存 data/processed/
\end{lstlisting}

這樣每次開啟專案，Claude Code 都會自動遵守這些規則，不需要重複說明。

\subsection{建立專案資料夾結構}

直接用對話指示 Claude Code 建立完整資料夾：

\begin{quote}
幫我建立一個合成控制研究的資料夾，包含：data/raw/、data/processed/、code/、output/figures/、output/tables/、notes/
\end{quote}

Claude Code 會自動執行 \texttt{mkdir -p} 建好所有資料夾。

\section{資料下載與整理}

\subsection{從 World Bank API 下載資料}

World Bank 提供免費 REST API，可以直接在 R 裡呼叫。以下範例下載多國的專利申請數：

\begin{lstlisting}
# 載入套件
library(tidyverse)
library(jsonlite)  # 解析 JSON 回傳值

# 定義下載函數
# 輸入：國家 ISO3 代碼、World Bank 指標代號
# 輸出：整理好的 tibble（country, year, value）
fetch_wb <- function(iso3, indicator) {
  url <- paste0(
    "https://api.worldbank.org/v2/country/", iso3,
    "/indicator/", indicator,
    "?date=2000:2023&format=json&per_page=1000"
  )
  raw <- fromJSON(url)
  raw[[2]] |>
    select(date, value) |>
    mutate(iso3 = iso3, year = as.integer(date)) |>
    select(iso3, year, value)
}

# 下載 7 個捐贈國的本國人 + 外國人專利申請數
countries <- c("JPN", "KOR", "HKG", "SGP", "THA", "PHL", "IDN")

patent_resd <- map_dfr(countries, ~fetch_wb(.x, "IP.PAT.RESD"))
patent_nres <- map_dfr(countries, ~fetch_wb(.x, "IP.PAT.NRES"))

# 計算總申請數（與台灣 TIPO 定義一致）
patent_total <- left_join(patent_resd, patent_nres,
                          by = c("iso3", "year"),
                          suffix = c("_resd", "_nres")) |>
  mutate(total = value_resd + value_nres) |>
  select(iso3, year, total)
\end{lstlisting}

\subsection{從 IMF API 下載台灣資料}

\begin{lstlisting}
# IMF WEO API — 指標 NGDPDPC = GDP per capita（現價美元）
fetch_imf <- function(indicator, iso3) {
  url <- paste0(
    "https://www.imf.org/external/datamapper/api/v1/",
    indicator, "/", iso3
  )
  raw <- fromJSON(url)
  values <- raw$values[[indicator]][[iso3]]
  tibble(
    iso3  = iso3,
    year  = as.integer(names(values)),
    value = as.numeric(values)
  )
}

# 下載台灣 GDP per capita
twn_gdp <- fetch_imf("NGDPDPC", "TWN") |>
  filter(year >= 2000, year <= 2023) |>
  rename(gdp_pc_usd = value)
\end{lstlisting}

\subsection{處理手動填入的資料}

讀取帶有 \texttt{\#} 注解行的 CSV：

\begin{lstlisting}
# comment = "#" 讓 readr 略過以 # 開頭的說明行
taiwan_raw <- read_csv("taiwan_patents_MANUAL.csv",
                       comment = "#",
                       show_col_types = FALSE)
\end{lstlisting}

\section{讓 Claude Code 除錯}

資料整理最常見的問題之一是路徑錯誤。例如 \texttt{here::here()} 找不到專案根目錄：

\begin{lstlisting}[language=bash]
Error: '/Users/xxx/Desktop/NTU_ECON/data/raw/patent_total.csv' does not exist.
\end{lstlisting}

直接把錯誤訊息貼給 Claude Code，它會自動診斷根本原因並修正：

\begin{lstlisting}[language=bash]
touch ~/Desktop/NTU_ECON/Course/LaborEcon/SSM_Taiwan_SC/.here
\end{lstlisting}

\section{實際效益}

\begin{center}
\begin{tabular}{lcc}
\toprule
\textbf{任務} & \textbf{沒有 AI} & \textbf{有 Claude Code} \\
\midrule
建資料夾結構 & 約 5 分鐘 & $<$ 30 秒 \\
寫 World Bank API 函數 & 約 30 分鐘 & 約 3 分鐘 \\
除錯路徑錯誤 & 約 15 分鐘 & $<$ 2 分鐘 \\
補齊台灣 GDP / R\&D 資料 & 約 1 小時 & 約 10 分鐘 \\
寫程式碼說明 & 約 20 分鐘 & $<$ 1 分鐘 \\
\bottomrule
\end{tabular}
\end{center}

\section{實用提示}

\begin{itemize}
  \item \textbf{給清楚的指令}：不要只說「幫我整理資料」，要說明具體的鍵值、欄位名稱和輸出格式。
  \item \textbf{讓 Claude Code 記住規則}：重要規則寫進 \texttt{CLAUDE.md}，不用每次重複說。
  \item \textbf{貼完整的錯誤訊息}：包含 traceback，Claude Code 才能定位問題。
  \item \textbf{分段確認}：複雜流程分段執行，每段確認結果正確後再繼續。
\end{itemize}

\section*{小結}

\begin{center}
\begin{tabular}{ll}
\toprule
\textbf{使用場景} & \textbf{Claude Code 的角色} \\
\midrule
建立專案結構 & 自動建資料夾、初始化 \texttt{.here}、\texttt{CLAUDE.md} \\
資料下載 & 撰寫 API 呼叫函數（World Bank、IMF） \\
資料整理 & 撰寫 \texttt{dplyr} pipeline，處理合併、清理 \\
除錯 & 診斷錯誤訊息，提出修正方案 \\
文件化 & 為每段程式碼加注解，撰寫資料來源文件 \\
\bottomrule
\end{tabular}
\end{center}

Claude Code 最大的優勢不是「取代你思考」，而是\textbf{大幅降低繁瑣實作的時間成本}，讓你可以把更多精力放在真正重要的事：研究設計、結果詮釋、論文撰寫。

\end{document}
